% ============================================================
%  CUADERNILLO DE ESTUDIO — CLASE 12
%  Seguridad de la Información y de Sistemas
%  Lic. en Gestión de Negocios Digitales — UCA
%  Compilar con: pdflatex (dos pasadas para el índice)
% ============================================================

\documentclass[12pt,a4paper]{article}

% ---------- Paquetes ----------
\usepackage[utf8]{inputenc}
\usepackage[spanish]{babel}
\usepackage[T1]{fontenc}
\usepackage{lmodern}
\usepackage[margin=2.2cm,headheight=14pt]{geometry}
\usepackage{amsmath}
\usepackage{amssymb}
\usepackage{enumitem}
\usepackage{booktabs}
\usepackage{array}
\usepackage{tabularx}
\usepackage[table]{xcolor}
\usepackage{tcolorbox}
\usepackage{graphicx}
\usepackage{fancyhdr}
\usepackage{titlesec}
\usepackage{hyperref}
\usepackage{multicol}
\usepackage{pifont}
\usepackage{wasysym}
\usepackage{tikz}
\usetikzlibrary{positioning,arrows.meta,shapes.geometric}

% ---------- Colores ----------
\definecolor{azulUCA}{HTML}{003366}
\definecolor{doradoUCA}{HTML}{B8860B}
\definecolor{grisClaro}{HTML}{F2F2F2}
\definecolor{rojoAlerta}{HTML}{C0392B}
\definecolor{verdeOK}{HTML}{27AE60}
\definecolor{azulCielo}{HTML}{2980B9}
\definecolor{naranjaAmenaza}{HTML}{D35400}
\definecolor{violetaIA}{HTML}{8E44AD}
\definecolor{rojoCrisis}{HTML}{922B21}
\definecolor{azulGobierno}{HTML}{1B4F72}

% Colores para los roles de la Junta de Crisis
\definecolor{rolGerente}{HTML}{D4AC0D}
\definecolor{rolLegal}{HTML}{2E4053}
\definecolor{rolComunicaciones}{HTML}{C0392B}
\definecolor{rolTI}{HTML}{1A5276}

% ---------- tcolorbox estilos ----------
\tcbuselibrary{skins,breakable}

\newtcolorbox{cajaTitulo}[1][]{
  colback=azulUCA, colframe=azulUCA, coltext=white,
  fonttitle=\Large\bfseries, arc=4mm, #1
}

\newtcolorbox{cajaDefinicion}{
  colback=azulCielo!10, colframe=azulCielo,
  fonttitle=\bfseries, title=Definici\'on, arc=3mm, breakable
}

\newtcolorbox{cajaCrisis}{
  colback=rojoCrisis!8, colframe=rojoCrisis,
  fonttitle=\bfseries\color{rojoCrisis},
  title={\ding{55}\; JUNTA DE CRISIS}, arc=3mm, breakable
}

\newtcolorbox{cajaInyeccion}{
  colback=naranjaAmenaza!12, colframe=naranjaAmenaza,
  fonttitle=\bfseries\color{naranjaAmenaza},
  title={\ding{115}\; Punto de inyecci\'on}, arc=3mm, breakable
}

\newtcolorbox{cajaActividad}{
  colback=grisClaro, colframe=azulUCA,
  fonttitle=\bfseries\color{azulUCA},
  title={\ding{45}\; Actividad Pr\'actica}, arc=3mm, breakable
}

\newtcolorbox{cajaNota}{
  colback=doradoUCA!10, colframe=doradoUCA,
  fonttitle=\bfseries, title=Concepto clave, arc=3mm, breakable
}

\newtcolorbox{cajaRolGerente}{
  colback=rolGerente!10, colframe=rolGerente,
  fonttitle=\bfseries\color{rolGerente},
  title={Rol: GERENCIA GENERAL}, arc=3mm, breakable
}

\newtcolorbox{cajaRolLegal}{
  colback=rolLegal!10, colframe=rolLegal,
  fonttitle=\bfseries\color{rolLegal},
  title={Rol: LEGAL / CUMPLIMIENTO}, arc=3mm, breakable
}

\newtcolorbox{cajaRolComunicaciones}{
  colback=rolComunicaciones!10, colframe=rolComunicaciones,
  fonttitle=\bfseries\color{rolComunicaciones},
  title={Rol: COMUNICACIONES}, arc=3mm, breakable
}

\newtcolorbox{cajaRolTI}{
  colback=rolTI!10, colframe=rolTI,
  fonttitle=\bfseries\color{rolTI},
  title={Rol: TI / SEGURIDAD}, arc=3mm, breakable
}

% ---------- Header / Footer ----------
\pagestyle{fancy}
\fancyhf{}
\fancyhead[L]{\small\textcolor{azulUCA}{Seguridad de la Informaci\'on y de Sistemas}}
\fancyhead[R]{\small\textcolor{azulUCA}{Cuadernillo --- Clase 12}}
\fancyfoot[C]{\thepage}
\fancyfoot[L]{\small\textcolor{gray}{UCA}}
\renewcommand{\headrulewidth}{0.4pt}
\renewcommand{\footrulewidth}{0.4pt}

% ---------- Títulos de sección ----------
\titleformat{\section}
  {\color{azulUCA}\Large\bfseries}
  {\thesection.}{0.6em}{}
\titleformat{\subsection}
  {\color{azulCielo}\large\bfseries}
  {\thesubsection}{0.6em}{}

% ---------- Enlaces ----------
\hypersetup{colorlinks=true, linkcolor=azulUCA, urlcolor=azulCielo}

% ---------- Checkbox ----------
\newcommand{\checkboxVacio}{$\square$\;}
\newcommand{\checkMark}{\ding{51}}

% ============================================================
\begin{document}

% ============================================================
%  PORTADA
% ============================================================
\begin{titlepage}
\centering
\vspace*{1cm}

\begin{cajaTitulo}[width=\textwidth, halign=center]
  SEGURIDAD DE LA INFORMACI\'ON Y DE SISTEMAS\\[4pt]
  {\large Licenciatura en Gesti\'on de Negocios Digitales --- UCA}
\end{cajaTitulo}

\vspace{1.5cm}

{\Huge\bfseries\color{azulUCA} CUADERNILLO DE ESTUDIO}\\[8pt]
{\LARGE\color{doradoUCA} CLASE 12}\\[12pt]
{\Large\itshape ``Factor humano y gesti\'on de incidentes:\\[2pt] Junta de Crisis''}\\[6pt]
{\large\color{rojoCrisis}\ding{55}\; Incluye simulaci\'on de filtraci\'on de datos}\\[1.5cm]

\begin{tabular}{rl}
\textbf{Profesor:}   & Mg.\ Ing.\ Rodrigo Atilio Elgueta \\[4pt]
\textbf{Unidad:}     & Unidad 4 \\[4pt]
\end{tabular}

\vfill

\begin{tcolorbox}[colback=grisClaro, colframe=azulUCA, arc=3mm, width=0.9\textwidth]
  \centering
  {\bfseries Datos del/la estudiante}\\[10pt]
  Nombre y apellido: \underline{\hspace{5cm}}\\[10pt]
  Grupo N.\textdegree: \underline{\hspace{5cm}}\\[10pt]
  Negocio Digital (TP): \underline{\hspace{5cm}}\\[10pt]
  Fecha: \underline{\hspace{5cm}}
\end{tcolorbox}

\vspace{1cm}
{\small\textcolor{gray}{Material de uso exclusivo para la cursada.}}
\end{titlepage}

% ============================================================
%  ÍNDICE
% ============================================================
\tableofcontents
\newpage

% ============================================================
%  SECCIÓN 1 — EL FACTOR HUMANO
% ============================================================
\section{El factor humano: el eslab\'on m\'as d\'ebil y el m\'as fuerte}

En la Clase 3 vimos que el \textbf{phishing}, el \textbf{ransomware} y la \textbf{ingenier\'ia social} son las amenazas m\'as comunes para los negocios digitales. En la Clase 10, practicamos c\'omo dise\~nar campa\~nas de concientizaci\'on. Hoy vamos a ver c\'omo \textbf{todo eso se pone a prueba} cuando las cosas salen mal.

\begin{cajaDefinicion}
El \textbf{factor humano} se refiere al rol de las personas dentro del ecosistema de seguridad de una organizaci\'on. Los empleados son simult\'aneamente:
\begin{itemize}[leftmargin=2em]
    \item \textbf{El eslab\'on m\'as d\'ebil:} un clic en un enlace de phishing puede comprometer toda la organizaci\'on.
    \item \textbf{La primera l\'inea de defensa:} un empleado bien capacitado puede detectar y reportar un ataque antes de que cause da\~no.
\end{itemize}
\end{cajaDefinicion}

\subsection{¿Por qu\'e las personas son cr\'iticas en seguridad?}

\renewcommand{\arraystretch}{1.5}
\begin{tabularx}{\textwidth}{|>{\bfseries\color{azulUCA}}p{3.5cm}|X|}
\hline
\rowcolor{grisClaro}
\textbf{Realidad} & \textbf{Implicancia} \\
\hline
95\% de los incidentes de seguridad tienen un componente humano (error, negligencia, manipulación) & Ninguna tecnología por sí sola puede proteger si las personas no saben cómo actuar. \\
\hline
Los atacantes prefieren el camino más fácil: engañar a una persona en lugar de hackear un sistema & La ingeniería social es más barata y efectiva que el hacking técnico. \\
\hline
Las políticas y controles existen, pero requieren que las personas las cumplan & Una PSA (Clase 11) no sirve si nadie la conoce o no hay capacitación. \\
\hline
\end{tabularx}

\begin{cajaNota}[title={La conexión con lo anterior}]
\begin{itemize}[leftmargin=2em]
    \item \textbf{Clase 10:} las campañas de concientización y los simulacros de phishing son \textbf{controles humanos} que reducen la probabilidad de incidentes.
    \item \textbf{Clase 11:} la PSA y el gobierno de la seguridad definen las \textbf{reglas del juego} que las personas deben seguir.
    \item \textbf{Clase 12 (hoy):} cuando todo lo anterior falla y ocurre un incidente real, la organizaci\'on necesita un \textbf{proceso de respuesta} y personas capacitadas para ejecutarlo.
\end{itemize}
\end{cajaNota}

\subsection{Programas de concientizaci\'on y capacitaci\'on}

Un programa efectivo de concientizaci\'on no es un evento \'unico (``una charla de una hora al a\~no''). Es un \textbf{proceso continuo} que incluye:

\begin{itemize}[leftmargin=2em]
    \item \textbf{Capacitaci\'on inicial} al ingreso (PSA, pol\'iticas, expectativas).
    \item \textbf{Actualizaciones peri\'odicas} (mensuales o trimestrales).
    \item \textbf{Simulacros de phishing} (Clase 10) para medir efectividad.
    \item \textbf{Micro-learning:} videos cortos, infograf\'ias, casos reales.
    \item \textbf{M\'etricas:} tasa de reporte de phishing, tiempo de respuesta, incidentes evitados.
\end{itemize}

\newpage
% ============================================================
%  SECCIÓN 2 — GESTIÓN DE INCIDENTES
% ============================================================
\section{Gesti\'on de incidentes de seguridad}

\begin{cajaDefinicion}
Un \textbf{incidente de seguridad} es cualquier evento que compromete la confidencialidad, integridad o disponibilidad de los activos de informaci\'on. La \textbf{gesti\'on de incidentes} es el proceso estructurado para detectarlos, contenerlos, erradicarlos, recuperar la normalidad y aprender de lo ocurrido.
\end{cajaDefinicion}

\subsection{El ciclo de vida de un incidente}

\begin{center}
\begin{tikzpicture}[
  node distance=0.8cm,
  box/.style={rectangle, draw=rojoCrisis, fill=rojoCrisis!10,
              text width=2.8cm, align=center, rounded corners=3pt,
              font=\small\bfseries},
  arr/.style={-{Stealth[length=3mm]}, thick, rojoCrisis}
]
  \node[box] (prep) {PREPARACI\'ON\\{\scriptsize Antes del incidente}};
  \node[box, right=1.5cm of prep] (det) {DETECCI\'ON\\{\scriptsize ¿C\'omo lo descubrimos?}};
  \node[box, right=1.5cm of det] (cont) {CONTENCI\'ON\\{\scriptsize Evitar que se expanda}};
  \node[box, right=1.5cm of cont] (err) {ERRADICACI\'ON\\{\scriptsize Eliminar la causa}};
  \node[box, below=1.5cm of err] (rec) {RECUPERACI\'ON\\{\scriptsize Volver a operar}};
  \node[box, left=1.5cm of rec] (lec) {LECCIONES\\{\scriptsize ¿Qu\'e aprendimos?}};

  \draw[arr] (prep) -- (det);
  \draw[arr] (det) -- (cont);
  \draw[arr] (cont) -- (err);
  \draw[arr] (err) -- (rec);
  \draw[arr] (rec) -- (lec);
  \draw[arr, dashed] (lec) -- (prep);
\end{tikzpicture}
\end{center}

\subsection{Las 4 fases en detalle}

\renewcommand{\arraystretch}{1.5}
\begin{tabularx}{\textwidth}{|>{\bfseries\color{rojoCrisis}}p{3cm}|X|X|}
\hline
\rowcolor{grisClaro}
\textbf{Fase} & \textbf{¿Qu\'e se hace?} & \textbf{¿Qui\'en lidera?} \\
\hline
\textbf{Detecci\'on} & Identificar el incidente, evaluar su alcance y severidad. & TI / Seguridad \\
\hline
\textbf{Respuesta} & Contener el da\~no, erradicar la amenaza, preservar evidencia. & TI / Seguridad + Legal \\
\hline
\textbf{Comunicaci\'on} & Informar a empleados, clientes, reguladores, medios. & Comunicaciones + Legal \\
\hline
\textbf{Recuperaci\'on} & Restaurar sistemas, monitorear, aprender del incidente. & Todos los roles \\
\hline
\end{tabularx}

\begin{cajaNota}[title={La diferencia entre un incidente y una crisis}]
Un \textbf{incidente} es un evento de seguridad que puede ser manejado por el equipo operativo. Una \textbf{crisis} es un incidente que \textbf{escala}: afecta la reputaci\'on, requiere comunicaci\'on externa, involucra a reguladores o tiene impacto financiero significativo. Hoy vamos a simular una \textbf{crisis}.
\end{cajaNota}

\newpage
% ============================================================
%  SECCIÓN 3 — LOS ROLES EN UNA CRISIS
% ============================================================
\section{Los 4 roles de la Junta de Crisis}

En la simulaci\'on de hoy, cada integrante del grupo asumir\'a uno de estos roles. Cada rol tiene responsabilidades espec\'ificas y debe tomar decisiones bajo presi\'on de tiempo.

\subsection{Resumen de roles}

\renewcommand{\arraystretch}{1.5}
\begin{tabularx}{\textwidth}{|>{\bfseries}p{3cm}|X|}
\hline
\rowcolor{grisClaro}
\textbf{Rol} & \textbf{Responsabilidad principal} \\
\hline
\textbf{Gerencia General} & Toma las decisiones estrat\'egicas finales. Aprueba presupuesto de emergencia. Decide el nivel de comunicaci\'on p\'ublica. \\
\hline
\textbf{Legal / Cumplimiento} & Asegura el cumplimiento normativo (Ley 25.326, GDPR). Eval\'ua responsabilidades legales. Coordina con abogados externos. \\
\hline
\textbf{Comunicaciones} & Gestiona la comunicaci\'on interna y externa. Redacta comunicados. Maneja redes sociales y prensa. \\
\hline
\textbf{TI / Seguridad} & Contiene t\'ecnicamente el incidente. Coordina la investigaci\'on forense. Restaura sistemas. \\
\hline
\end{tabularx}

\newpage
\subsection{Tarjetas de rol}

Cada estudiante debe recibir (o recortar) una tarjeta de rol antes de comenzar la simulaci\'on.

\begin{cajaRolGerente}
\textbf{Tu misi\'on:} Proteger el negocio a largo plazo.\\[4pt]
\textbf{Tus decisiones clave:}
\begin{itemize}[leftmargin=1.5em]
    \item ¿Cu\'anto presupuesto de emergencia aprob\'as para contener el incidente?
    \item ¿Notificamos a los clientes ahora o esperamos a tener m\'as informaci\'on?
    \item ¿Hacemos declaraci\'on p\'ublica o esperamos a que la prensa nos contacte?
    \item ¿Contratamos un equipo forense externo?
\end{itemize}
\textbf{Tu pregunta recurrente:} ``¿Cu\'al es el impacto financiero y reputacional de cada opci\'on?''
\end{cajaRolGerente}

\vspace{6pt}

\begin{cajaRolLegal}
\textbf{Tu misi\'on:} Minimizar la responsabilidad legal y cumplir la ley.\\[4pt]
\textbf{Tus decisiones clave:}
\begin{itemize}[leftmargin=1.5em]
    \item ¿Cu\'ando hay que notificar a la AAIP (Agencia de Acceso a la Informaci\'on P\'ublica)?
    \item ¿La Ley 25.326 exige notificaci\'on a los titulares de los datos?
    \item ¿Hay clientes en Europa? ¿Aplica el GDPR?
    \item ¿Qu\'e evidencias hay que preservar para una eventual investigaci\'on judicial?
\end{itemize}
\textbf{Tu pregunta recurrente:} ``¿Qu\'e dice la ley y cu\'ales son los plazos?''
\end{cajaRolLegal}

\vspace{6pt}

\begin{cajaRolComunicaciones}
\textbf{Tu misi\'on:} Proteger la reputaci\'on y la confianza de los clientes.\\[4pt]
\textbf{Tus decisiones clave:}
\begin{itemize}[leftmargin=1.5em]
    \item ¿Qu\'e decimos en el comunicado interno a los empleados?
    \item ¿Qu\'e decimos en el comunicado externo a los clientes?
    \item ¿C\'omo respondemos si la prensa nos contacta?
    \item ¿Publicamos algo en redes sociales antes de que se viralice?
\end{itemize}
\textbf{Tu pregunta recurrente:} ``¿Qu\'e van a pensar los clientes y c\'omo van a reaccionar?''
\end{cajaRolComunicaciones}

\vspace{6pt}

\begin{cajaRolTI}
\textbf{Tu misi\'on:} Contener t\'ecnicamente el incidente y restaurar la operaci\'on.\\[4pt]
\textbf{Tus decisiones clave:}
\begin{itemize}[leftmargin=1.5em]
    \item ¿Desconectamos el servidor afectado de la red?
    \item ¿Apagamos los sistemas comprometidos o los aislamos?
    \item ¿Preservamos evidencia forense (logs, im\'agenes de disco)?
    \item ¿Cu\'anto tiempo estimamos para restaurar la operaci\'on normal?
\end{itemize}
\textbf{Tu pregunta recurrente:} ``¿Cu\'al es el alcance t\'ecnico del incidente y c\'omo evitamos que se expanda?''
\end{cajaRolTI}

\newpage
% ============================================================
%  SECCIÓN 4 — LA SIMULACIÓN: JUNTA DE CRISIS
% ============================================================
\section{La simulaci\'on: Junta de Crisis}

\begin{cajaCrisis}[title={Instrucciones de la simulaci\'on}]
\textbf{Duraci\'on total:} 75 minutos de simulaci\'on + 15 minutos de debrief.\\[4pt]
\textbf{Mec\'anica:}
\begin{enumerate}[leftmargin=2em]
    \item El docente presenta el escenario inicial (la ``bomba'' cae).
    \item Cada grupo se re\'une como Junta de Crisis y trabaja en tiempo real.
    \item El docente introduce \textbf{puntos de inyecci\'on} (eventos sorpresa) en momentos espec\'ificos para forzar decisiones.
    \item Cada grupo completa las plantillas de respuesta de su rol.
    \item Al final, cada grupo presenta sus decisiones y se hace el debrief.
\end{enumerate}
\end{cajaCrisis}

\subsection{Escenario inicial: ``DataShop''}

\begin{tcolorbox}[colback=grisClaro, colframe=rojoCrisis, arc=3mm]
\textbf{Contexto:} ``DataShop'' es un e-commerce de electr\'onica con programa de fidelidad. Tiene 8.000 clientes activos y una base de datos con nombres, emails, tel\'efonos e historial de compras.\\[4pt]
\textbf{Lunes 9:15 AM:} El equipo de TI detecta tr\'afico de salida inusual desde el servidor de la base de datos.\\[4pt]
\textbf{Investigaci\'on inicial (9:15-9:30):}
\begin{itemize}[leftmargin=1.5em]
    \item La base de datos de clientes fue accedida desde una IP externa.
    \item Se encontr\'o un archivo ejecutable sospechoso en la PC de un empleado del \'area de Finanzas.
    \item El empleado recuerda haber abierto un adjunto en un email que parec\'ia ser de un proveedor.
    \item El adjunto instal\'o software de acceso remoto.
\end{itemize}
\textbf{Datos potencialmente comprometidos:}
\begin{itemize}[leftmargin=1.5em]
    \item Nombres, emails, tel\'efonos, historial de compras (8.000 clientes).
    \item Datos de pago (cifrados, pero no se sabe si la clave de cifrado fue comprometida).
\end{itemize}
\textbf{Situaci\'on adicional:} Un foro de hackers public\'o una muestra de datos de DataShop y el atacante est\'a exigiendo un rescate.
\end{tcolorbox}

\newpage
\subsection{Cronograma de la simulaci\'on y puntos de inyecci\'on}

La simulaci\'on se desarrolla en tiempo real. El docente introduce \textbf{puntos de inyecci\'on} (eventos sorpresa) en los momentos indicados para forzar decisiones.

\renewcommand{\arraystretch}{1.6}
\begin{tabularx}{\textwidth}{|>{\bfseries\color{rojoCrisis}}p{2cm}|X|}
\hline
\rowcolor{grisClaro}
\textbf{Tiempo} & \textbf{Evento / Punto de inyecci\'on} \\
\hline
T = 0 min & Se presenta el escenario inicial. La Junta de Crisis se re\'une. \\
\hline
T = 10 min & \textbf{Inyecci\'on 1:} Un periodista env\'ia un email preguntando si DataShop sufri\'o una filtraci\'on de datos. \\
\hline
T = 25 min & \textbf{Inyecci\'on 2:} Un cliente publica en redes sociales que recibi\'o un email sospechoso a nombre de DataShop. \\
\hline
T = 40 min & \textbf{Inyecci\'on 3:} El atacante env\'ia una nota de rescate con un plazo de 48 horas para pagar. \\
\hline
T = 55 min & \textbf{Inyecci\'on 4:} La AAIP env\'ia una notificaci\'on solicitando informaci\'on sobre una posible brecha de datos. \\
\hline
T = 70 min & Cierre de la simulaci\'on. Cada grupo presenta sus decisiones finales. \\
\hline
T = 75-90 min & Debrief grupal: ¿qu\'e funcion\'o, qu\'e no, qu\'e aprendimos? \\
\hline
\end{tabularx}

\newpage
% ============================================================
%  SECCIÓN 5 — PLANTILLAS DE RESPUESTA
% ============================================================
\section{Plantillas de respuesta por rol}

Durante la simulaci\'on, cada rol debe completar estas plantillas. Al final, se presentan al docente y al resto de la clase.

\subsection{Plantilla del Gerente General}

\begin{cajaRolGerente}
\renewcommand{\arraystretch}{1.6}
\begin{tabularx}{\textwidth}{|X|}
\hline
\textbf{Decisi\'on 1: Presupuesto de emergencia}\\[2pt]
?Cu\'anto presupuesto aprueban para contener el incidente? (forense externo, consultor\'ia legal, etc.)\\[1.5cm]
\hline
\textbf{Decisi\'on 2: Notificaci\'on a clientes}\\[2pt]
?Notificamos a los 8.000 clientes ahora o esperamos? Justificaci\'on:\\[1.5cm]
\hline
\textbf{Decisi\'on 3: Comunicaci\'on p\'ublica}\\[2pt]
?Hacemos declaraci\'on p\'ublica proactiva o esperamos a que la prensa nos contacte? Justificaci\'on:\\[1.5cm]
\hline
\textbf{Decisi\'on 4: Rescate}\\[2pt]
?Pagamos el rescate al atacante? ?Por qu\'e s\'i o por qu\'e no? (Considerar implicancias legales y \'eticas):\\[1.5cm]
\hline
\end{tabularx}
\end{cajaRolGerente}

\newpage
\subsection{Plantilla del Legal / Cumplimiento}

\begin{cajaRolLegal}
\renewcommand{\arraystretch}{1.6}
\begin{tabularx}{\textwidth}{|X|}
\hline
\textbf{Cuesti\'on 1: Obligaci\'on de notificaci\'on a la AAIP}\\[2pt]
?Qu\'e dice la Ley 25.326 sobre la notificaci\'on de brechas de datos? ?Cu\'al es el plazo?\\[1.5cm]
\hline
\textbf{Cuesti\'on 2: Notificaci\'on a los titulares de datos}\\[2pt]
?Los 8.000 clientes deben ser notificados? ?Qu\'e tipo de datos se vieron comprometidos?\\[1.5cm]
\hline
\textbf{Cuesti\'on 3: Aplicabilidad del GDPR}\\[2pt]
?Hay clientes en Europa? ?Aplica el GDPR? ?Hay obligaciones adicionales?\\[1.5cm]
\hline
\textbf{Cuesti\'on 4: Preservaci\'on de evidencia}\\[2pt]
?Qu\'e evidencias deben preservarse para una eventual investigaci\'on judicial?\\[1.5cm]
\hline
\end{tabularx}
\end{cajaRolLegal}

\newpage
\subsection{Plantilla de Comunicaciones}

\begin{cajaRolComunicaciones}
\renewcommand{\arraystretch}{1.6}
\begin{tabularx}{\textwidth}{|X|}
\hline
\textbf{Comunicado interno a empleados}\\[2pt]
Redacten el comunicado interno (2-3 p\'arrafos):\\[3cm]
\hline
\textbf{Comunicado externo a clientes}\\[2pt]
Redacten el comunicado externo (2-3 p\'arrafos):\\[3cm]
\hline
\textbf{Respuesta a prensa}\\[2pt]
Redacten una respuesta de 2-3 oraciones para el periodista que envi\'o el email:\\[2cm]
\hline
\textbf{Redes sociales}\\[2pt]
?Publican algo en redes sociales antes de que se viralice? ?Qu\'e y cu\'ando?\\[1.5cm]
\hline
\end{tabularx}
\end{cajaRolComunicaciones}

\newpage
\subsection{Plantilla de TI / Seguridad}

\begin{cajaRolTI}
\renewcommand{\arraystretch}{1.6}
\begin{tabularx}{\textwidth}{|X|}
\hline
\textbf{Acci\'on 1: Contenci\'on inmediata}\\[2pt]
?Desconectan el servidor afectado de la red? ?Apagan los sistemas comprometidos? Justifiquen:\\[2cm]
\hline
\textbf{Acci\'on 2: Preservaci\'on de evidencia}\\[2pt]
?Qu\'e evidencias t\'ecnicas deben preservar? (logs, im\'agenes de disco, capturas de red):\\[2cm]
\hline
\textbf{Acci\'on 3: An\'alisis forense}\\[2pt]
?Contratan un equipo forense externo? ?Qu\'e esperan que determinen?\\[2cm]
\hline
\textbf{Acci\'on 4: Restauraci\'on}\\[2pt]
?Cu\'anto tiempo estiman para restaurar la operaci\'on normal? ?Qu\'e sistemas se restauran primero?\\[2cm]
\hline
\end{tabularx}
\end{cajaRolTI}

\newpage
% ============================================================
%  SECCIÓN 6 — DEBRIEF Y LECCIONES
% ============================================================
\section{Debrief y lecciones aprendidas}

Al final de la simulaci\'on, cada grupo presenta sus decisiones y se hace un debrief grupal. Este es el momento de reflexionar sobre lo que funcion\'o, lo que no, y qu\'e habr\'ian hecho diferente.

\begin{cajaActividad}[title={Reflexi\'on grupal}]
Respondan como grupo:
\end{cajaActividad}

\vspace{6pt}
\renewcommand{\arraystretch}{1.6}
\begin{tabularx}{\textwidth}{|X|}
\hline
\textbf{1. ?Qu\'e decisi\'on fue la m\'as dif\'icil de tomar? ?Por qu\'e?}\\[2.5cm]
\hline
\textbf{2. ?Qu\'e rol tuvo m\'as presi\'on durante la simulaci\'on? ?C\'omo lo manejaron?}\\[2.5cm]
\hline
\textbf{3. ?Qu\'e habr\'ian hecho diferente si hubieran tenido m\'as tiempo o m\'as informaci\'on?}\\[2.5cm]
\hline
\textbf{4. ?Qu\'e elemento de la PSA (Clase 11) habr\'ia ayudado a prevenir o manejar mejor esta crisis?}\\[2.5cm]
\hline
\textbf{5. ?Qu\'e incluir\'ian en su Plan de Respuesta a Incidentes del TP Integrador a partir de esta experiencia?}\\[3cm]
\hline
\end{tabularx}

\begin{cajaNota}[title={Conexi\'on con el TP Integrador}]
Lo que aprendan hoy en la Junta de Crisis es el \textbf{insumo directo} para el Entregable 6 del TP Integrador: \textbf{Plan de respuesta a incidentes y de continuidad de negocio}. Cada decisi\'on que tomaron hoy es una pieza de ese plan. Guarden las plantillas completadas.
\end{cajaNota}

\newpage
% ============================================================
%  SECCIÓN 7 — GLOSARIO Y NOTAS
% ============================================================
\section{Glosario de la Clase 12}

\renewcommand{\arraystretch}{1.4}
\begin{tabularx}{\textwidth}{>{\bfseries}p{4.5cm}X}
\toprule
\textbf{T\'ermino} & \textbf{Definici\'on en una l\'inea} \\
\midrule
Incidente de seguridad & Evento que compromete la CID de los activos de informaci\'on. \\
Crisis & Incidente que escala y requiere comunicaci\'on externa o tiene impacto significativo. \\
Junta de Crisis & Equipo multidisciplinario que gestiona la respuesta a un incidente grave. \\
Detecci\'on & Fase de identificaci\'on del incidente. \\
Contenci\'on & Fase de limitar el da\~no y evitar que el incidente se expanda. \\
Erradicaci\'on & Fase de eliminar la causa ra\'iz del incidente. \\
Recuperaci\'on & Fase de restaurar la operaci\'on normal. \\
Lecciones aprendidas & An\'alisis posterior al incidente para mejorar futuros procesos. \\
Concientizaci\'on & Capacitaci\'on para que las personas reconozcan y eviten amenazas. \\
Inyecci\'on & Evento sorpresa introducido durante una simulaci\'on para forzar decisiones. \\
\bottomrule
\end{tabularx}

\vspace{1cm}

\section{Espacio para notas y preparaci\'on}

\begin{tcolorbox}[colback=grisClaro, colframe=gray!50, arc=2mm, height=7cm]
\textcolor{gray}{\itshape Us\'a este espacio para anotar decisiones clave de la simulaci\'on, ideas para tu Plan de Respuesta a Incidentes del TP, o dudas sobre la gesti\'on de crisis.}
\end{tcolorbox}

\vfill
\begin{center}
\small\textcolor{gray}{%
Cuadernillo elaborado para la c\'atedra de Seguridad de la Informaci\'on y de
Sistemas --- Lic.\ en Gesti\'on de Negocios Digitales, UCA.\\
Prohibida su distribuci\'on fuera del \'ambito de la cursada sin autorizaci\'on del docente.}
\end{center}

\end{document}